\documentclass[10pt,a4paper]{amsart}
\usepackage[margin=19mm]{geometry}
\usepackage{amsmath,amssymb,mathtools}
\usepackage[scr=rsfs,cal=euler]{mathalfa}
\usepackage{xfrac}
\usepackage[unicode,hidelinks,bookmarks=false]{hyperref}
\newcommand{\ie}{i.e.\ }
\newcommand{\cf}{cf.\ }
\newcommand{\resp}{respectively\ }
\newcommand{\Subsection}{\subsection}
\providecommand{\nobreakdash}{\nobreak\hbox{-}}
\setlength{\emergencystretch}{2em}
\multlinegap0pt
\usepackage{amssymb}
\usepackage[shortlabels]{enumitem}
\newtheorem{theorem}{Theorem}
\title{Drinfeld Correspondence in Infinite Dimensions}
\author{Praful Rahangdale}
\date{Source: arXiv:2603.04634v1 (CC BY 4.0)}
\begin{document}
\maketitle

\noindent\emph{Source and licence.} Drinfeld Correspondence in Infinite Dimensions. Version arXiv:2603.04634v1 (CC BY 4.0), distributed by arXiv under the Creative Commons Attribution 4.0 International licence, http://creativecommons.org/licenses/by/4.0/, as recorded in the arXiv licence metadata for this exact version and shown on the abstract page https://arxiv.org/abs/2603.04634v1. Full licence notice: \texttt{licenses/arxiv-2603.04634-CC-BY-4.0.txt}.\par\smallskip

\noindent\textit{Editorial scope.} Counted unit: the article's theorem theorem (source0.tex lines 1015--1032). The article's complete original proof of it is delivered with it (source0.tex lines 1034--1088, one \texttt{proof} environment of 55 lines). No mathematics is written, repaired or re-derived: the statement and the proof are the article's own lines, in the article's own notation, and no retained line is substituted or re-flowed.  No further source-local context is needed: every cross-reference of every retained line resolves inside the two delivered spans.  Nothing was omitted from any retained span, so the package carries no omission record.  No retained line contains a citation, so the wrapper carries no bibliography and no bibliography entry.  No retained line contains a figure environment, an image-inclusion command or drawing code, so no image is delivered and the package declares no asset dependency.  Editorial formatting only, in addition: an AMS \texttt{amsart} preamble that restates the article's own declarations and macros used by the retained lines, namely the environments \texttt{theorem} on separate counters, together with the standard packages those lines need.  The retained environments print in this document as Theorem 1 in order of appearance, whereas the article prints the same items with the numbering of its own full document.  The complete native source of the article as distributed in the arXiv e-print for this exact version is bundled unchanged as \texttt{sources/195752/source0.tex}.
\begin{theorem}
Let $A$ be a unital $C^*$-algebra equipped with a strongly continuous circle action $\alpha: \mathbb{S}^1 \rightarrow \text{Aut}(A)$ and a faithful, tracial, $\alpha$-invariant state $\omega: A \rightarrow \mathbb{C}$. Let $\pi_0: A_{\infty} \rightarrow A_{\infty}^0$ be the projection onto the zero-weight elements, given by
\[
\pi_0(a) := \frac{1}{2\pi} \int_0^{2\pi} \alpha_t(a) \, dt.
\]
Consider the Lie algebra $\mathfrak{g} := A_{\infty} \oplus A_{\infty}$ equipped with the componentwise Lie bracket $[(x_1, x_2), (y_1, y_2)] := ([x_1, y_1], [x_2, y_2])$ and the symmetric bilinear form
\[
\langle\langle (a, a'), (x, y) \rangle\rangle := \omega(ax) - \omega(a'y).
\]
Define the Lie subalgebras $\mathfrak{p}_1, \mathfrak{p}_2 \subset \mathfrak{g}$ by
\begin{align*}
\mathfrak{p}_1 &:= \{(a, a) \mid a \in A_{\infty}\}, \\
\mathfrak{p}_2 &:= \{(x, y) \mid x \in A_{\infty}^{\ge 0}, \, y \in A_{\infty}^{\le 0}, \, \pi_0(x) + \pi_0(y) = 0\},
\end{align*}
where $A_{\infty}^{\ge 0} := \bigoplus_{m \ge 0} A_{\infty}^m$ and $A_{\infty}^{\le 0} := \bigoplus_{m \le 0} A_{\infty}^m$. 

If the zero-weight subspace $A_{\infty}^0$ is abelian with respect to the canonical commutator bracket (i.e., $[a, b] = 0$ for all $a, b \in A_{\infty}^0$), then $(\mathfrak{g}, \mathfrak{p}_1, \mathfrak{p}_2)$ is a Manin triple.
\end{theorem}

\begin{proof}
We must verify that the triple $(\mathfrak{g}, \mathfrak{p}_1, \mathfrak{p}_2)$ satisfies the four defining conditions of a Manin triple:

(1) \textit{Lie subalgebras}. For $\mathfrak{p}_1 = \{(a, a) \mid a \in A_{\infty}\}$, the componentwise bracket yields
\[
[(a, a), (b, b)] = ([a, b], [a, b]) \in \mathfrak{p}_1,
\]
thus $\mathfrak{p}_1$ is closed under the Lie bracket. For $\mathfrak{p}_2$, let $(x_1, y_1), (x_2, y_2) \in \mathfrak{p}_2$. The bracket is
\[
[(x_1, y_1), (x_2, y_2)] = ([x_1, x_2], [y_1, y_2]).
\]
Writing $x_1 = \sum_{m \ge 0} x_{1,m}$ and $x_2 = \sum_{n \ge 0} x_{2,n}$ with elements in the respective weight spaces $A_{\infty}^m$ and $A_{\infty}^n$, we have
\[
[x_1, x_2] = \sum_{m, n \ge 0} [x_{1,m}, x_{2,n}].
\]
Because the circle action acts via $\alpha_t([x_{1,m}, x_{2,n}]) = e^{i(m+n)t}[x_{1,m}, x_{2,n}]$ and $m+n \ge 0$, we conclude that $[x_1, x_2] \in A_{\infty}^{\ge 0}$. By identical logic, $[y_1, y_2] \in A_{\infty}^{\le 0}$. Furthermore, because $[A_{\infty}^0, A_{\infty}^0] = 0$, the projection of the bracket onto the zero-weight space vanishes: $\pi_0([x_1, x_2]) = [\pi_0(x_1), \pi_0(x_2)] = 0$. Hence, $\pi_0([x_1, x_2]) + \pi_0([y_1, y_2]) = 0 + 0 = 0$, proving $\mathfrak{p}_2$ is a closed Lie subalgebra.

(2) \textit{Isotropy}. For any two elements $(a, a), (b, b) \in \mathfrak{p}_1$, the bilinear form gives
\[
\langle\langle (a, a), (b, b) \rangle\rangle = \omega(ab) - \omega(ab) = 0.
\]
For $\mathfrak{p}_2$, let $(x, y), (x', y') \in \mathfrak{p}_2$. Then $x, x' \in A_{\infty}^{\ge 0}$ and $y, y' \in A_{\infty}^{\le 0}$, and we evaluate
\[
\langle\langle (x, y), (x', y') \rangle\rangle = \omega(xx') - \omega(yy').
\]
Because the state $\omega$ is $\alpha$-invariant, $\omega(a) = \omega(\pi_0(a))$ for all $a \in A_\infty$). Since $x$ and $x'$ only contain weights $m, n \ge 0$, their product $xx'$ only has a zero-weight component if $m=0$ and $n=0$. Therefore, $\pi_0(xx') = \pi_0(x)\pi_0(x')$, which implies $\omega(xx') = \omega(\pi_0(x)\pi_0(x'))$. Similarly, $\omega(yy') = \omega(\pi_0(y)\pi_0(y'))$. Since $\pi_0(x) = -\pi_0(y)$ and $\pi_0(x') = -\pi_0(y')$, we obtain
\[
\omega(\pi_0(x)\pi_0(x')) - \omega(\pi_0(y)\pi_0(y')) = \omega((-\pi_0(y))(-\pi_0(y'))) - \omega(\pi_0(y)\pi_0(y')) = 0.
\]
Therefore, both $\mathfrak{p}_1$ and $\mathfrak{p}_2$ are isotropic.

(3) \textit{Invariance under the adjoint action}. The bilinear form is invariant if for all $(u, v), (a, b), (x, y) \in \mathfrak{g}$, we have
\[
\langle\langle [(u, v), (a, b)], (x, y) \rangle\rangle + \langle\langle (a, b), [(u, v), (x, y)] \rangle\rangle = 0.
\]
Expanding the brackets and the bilinear form yields
\[
\omega([u, a]x) - \omega([v, b]y) + \omega(a[u, x]) - \omega(b[v, y]).
\]
Because the state $\omega$ is tracial, we can cyclically permute the elements: $\omega([u,a]x) = \omega(uax) - \omega(aux)$ and $\omega(a[u,x]) = \omega(aux) - \omega(axu) = \omega(aux) - \omega(uax)$. Summing these two terms gives zero. The same cancellation occurs also for the $v, b, y$ terms, proving invariance.

(4) \textit{Weak non-degeneracy and topological direct sum}. Weak non-degeneracy of the pairing follows directly from the faithfulness of the state $\omega$. Indeed, if $\omega(ax) - \omega(by) = 0$ for all $(x, y) \in \mathfrak{g}$, setting $y=0$ and $x=a^*$ gives $\omega(aa^*) = 0$, which implies $a=0$ (and symmetrically $b=0$). 

For the topological direct sum, we must show that any element $(u, v) \in A_{\infty} \oplus A_{\infty}$ decomposes uniquely into $(a, a) \in \mathfrak{p}_1$ and $(x, y) \in \mathfrak{p}_2$. Let $P_+ = \sum_{m > 0}$ and $P_- = \sum_{m < 0}$ denote the projections onto the strictly positive and strictly negative weight spaces. The unique decomposition is given explicitly by
\[
a = P_+ v + \frac{1}{2}\pi_0(u+v) + P_- u,
\]
\[
x = P_+ (u-v) + \frac{1}{2}\pi_0(u-v),
\]
\[
y = \frac{1}{2}\pi_0(v-u) + P_- (v-u).
\]
Note that, $x \in A_{\infty}^{\ge 0}$ and $y \in A_{\infty}^{\le 0}$. Moreover, $\pi_0(x) + \pi_0(y) = \frac{1}{2}\pi_0(u-v) + \frac{1}{2}\pi_0(v-u) = 0$, so $(x,y) \in \mathfrak{p}_2$. Furthermore, $a + x = u$ and $a + y = v$. Both $\mathfrak{p}_1$ and $\mathfrak{p}_2$ are closed subspaces because the projections $P_+, P_-,$ and $\pi_0$ are continuous.
\end{proof}

\end{document}
